\chapter{Spectral sequences}
\label{app:spectral-sequences}

\index{spectral sequence|textbf}

Three spectral sequences are central to chiral Koszul duality: the bar spectral sequence
(computing bar complex homology), the Chevalley--Cousin spectral sequence
(relating D-module cohomology to configuration space geometry), and the genus
spectral sequence (organizing quantum corrections by loop order).

\section{Filtered complexes and spectral sequences}
\label{sec:filtered-complexes}

\begin{definition}[Filtered chain complex]
\label{def:filtered-complex}
\index{filtered complex}
A \emph{filtered chain complex} is a chain complex $(C, d)$ together with a 
sequence of subcomplexes:
\[
\cdots \subseteq F_p C \subseteq F_{p+1} C \subseteq \cdots \subseteq C
\]
such that $d(F_p C) \subseteq F_p C$ for all $p$.

The filtration is:
\begin{enumerate}[label=(\roman*)]
\item \emph{Exhaustive} if $C = \bigcup_p F_p C$.
\item \emph{Complete} (or Hausdorff) if $\bigcap_p F_p C = 0$.
\item \emph{Bounded below} if for each degree $n$, there exists $p_0(n)$
such that $F_p C_n = 0$ for $p < p_0(n)$.
\item \emph{Bounded above} if for each degree $n$, there exists $p_1(n)$
such that $F_p C_n = C_n$ for $p > p_1(n)$.
\item \emph{Bounded} if both bounded below and bounded above.
\end{enumerate}
\end{definition}

\begin{construction}[Associated graded and spectral sequence]
\label{constr:ass-graded-ss}
Given a filtered complex $(C, F_\bullet, d)$:

\emph{Step~1.} The \emph{associated graded} is:
\[
\gr_p C := F_p C / F_{p-1} C, \quad \gr C := \bigoplus_p \gr_p C.
\]
The differential $d$ induces $d_0: \gr_p C_n \to \gr_p C_{n-1}$ since 
$d(F_p C) \subseteq F_p C$.

\emph{Step~2.} Define the \emph{$r$-th approximation}:
\[
Z_r^{p,q} := \{x \in F_p C_{p+q} : dx \in F_{p-r} C_{p+q-1}\} / F_{p-1} C_{p+q}
\]
\[
B_r^{p,q} := \{dx : x \in F_{p+r-1} C_{p+q+1}, dx \in F_p C_{p+q}\} / F_{p-1} C_{p+q}
\]

\emph{Step~3.} The \emph{$E_r$-page} is:
\[
E_r^{p,q} := Z_r^{p,q} / B_r^{p,q}.
\]
The differential $d_r: E_r^{p,q} \to E_r^{p-r, q+r-1}$ is induced by $d$:
\[
d_r([x]) := [dx] \in E_r^{p-r, q+r-1}
\]
for $[x] \in E_r^{p,q}$ represented by $x \in F_p C_{p+q}$ with $dx \in F_{p-r} C_{p+q-1}$.

\begin{remark}[Convention]
This appendix presents spectral sequences in the \emph{homological} convention (chain complexes, $d$ lowers total degree), so $d_r$ has bidegree $(-r, r-1)$. The main body of the text uses cohomological convention ($d$ raises degree), where the differential $d_r: E_r^{p,q} \to E_r^{p+r, q-r+1}$ has bidegree $(r, -r+1)$. The two conventions are related by reindexing $E_r^{p,q} \leftrightarrow E_r^{-p,-q}$.
\end{remark}

\emph{Step~4.} There are canonical isomorphisms:
\[
H(E_r^{p,q}, d_r) \cong E_{r+1}^{p,q}.
\]
The sequence $(E_r, d_r)_{r \geq 0}$ is the \emph{spectral sequence} associated 
to the filtration.
\end{construction}

\begin{theorem}[Identification of early pages; \ClaimStatusProvedElsewhere]
\label{thm:early-pages}
For the spectral sequence of a filtered complex:
\begin{enumerate}[label=(\roman*)]
\item $E_0^{p,q} = \gr_p C_{p+q} = F_p C_{p+q} / F_{p-1} C_{p+q}$.
\item $d_0: E_0^{p,q} \to E_0^{p, q-1}$ is the induced differential on $\gr C$.
\item $E_1^{p,q} = H_{p+q}(\gr_p C) = H_{p+q}(F_p C / F_{p-1} C)$.
\item $d_1: E_1^{p,q} \to E_1^{p-1, q}$ is the \emph{connecting homomorphism} 
in the long exact sequence of the short exact sequence:
\[
0 \to \gr_{p-1} C \to F_p C / F_{p-2} C \to \gr_p C \to 0.
\]
\end{enumerate}
\end{theorem}

\begin{proof}
For (i) and (ii): By definition, $Z_0^{p,q} = F_p C_{p+q} / F_{p-1} C_{p+q}$ 
(all elements, since $F_{p-0} = F_p$ contains $d(F_p C_{p+q+1})$), and 
$B_0^{p,q} = 0$ (no elements of the form $dx$ with $x \in F_{p-1}$). Thus 
$E_0 = \gr C$.

For (iii): $E_1^{p,q} = Z_1^{p,q} / B_1^{p,q}$ where $Z_1^{p,q}$ consists of 
elements $[x]$ with $dx \in F_{p-1}$, i.e., $d_0[x] = 0$ in $\gr_p$. The 
quotient by $B_1^{p,q} = \im(d_0)$ gives $\ker(d_0)/\im(d_0) = H(\gr_p C)$.

For (iv): The differential $d_1$ maps an element $[x] \in H(\gr_p C)$ to $[dx] 
\in H(\gr_{p-1} C)$. This is the connecting map $\delta$ in the 
long exact sequence.
\end{proof}

\begin{definition}[Convergence]
\label{def:ss-convergence}
\index{spectral sequence!convergence}
A spectral sequence $(E_r^{p,q}, d_r)$ \emph{converges} to a graded object 
$H^* = \bigoplus_n H^n$ if there is a filtration $F_p H^n$ and isomorphisms:
\[
E_\infty^{p,q} \cong \gr_p H^{p+q} := F_p H^{p+q} / F_{p-1} H^{p+q}
\]
for all $p, q$, where:
\[
E_\infty^{p,q} := \bigcap_{r \geq 0} Z_r^{p,q} / \bigcup_{r \geq 0} B_r^{p,q}.
\]
We write $E_r^{p,q} \Rightarrow H^{p+q}$.
\end{definition}

\begin{theorem}[Classical convergence theorem \cite{Weibel94}; \ClaimStatusProvedElsewhere]
\label{thm:classical-convergence}
Let $(C, F_\bullet, d)$ be a filtered chain complex.
\begin{enumerate}[label=(\roman*)]
\item If the filtration is bounded, the spectral sequence converges to $H_*(C)$:
\[
E_r^{p,q} \Rightarrow H_{p+q}(C).
\]
The spectral sequence is \emph{regular}: $E_r = E_\infty$ for
$r$ sufficiently large (depending on $p, q$).

\item If the filtration is exhaustive and complete, and bounded below in each 
degree, then the spectral sequence converges conditionally:
\[
E_r^{p,q} \Rightarrow H_{p+q}(\widehat{C})
\]
where $\widehat{C} = \varprojlim_p C / F_p C$ is the completion.
\end{enumerate}
\end{theorem}


\section{Convergence criteria}
\label{sec:convergence-criteria}

Refined convergence criteria for pro-nilpotent and completed settings follow.

\begin{definition}[Regular spectral sequence]
\label{def:regular-ss}
A spectral sequence is \emph{regular} if for each bidegree $(p, q)$, the 
sequence stabilizes: there exists $r_0 = r_0(p, q)$ such that:
\[
E_r^{p,q} = E_{r+1}^{p,q} = \cdots = E_\infty^{p,q} \quad \text{for } r \geq r_0.
\]
Equivalently, $d_r: E_r^{p,q} \to E_r^{p-r, q+r-1}$ is zero for $r \geq r_0$, 
and $d_r: E_r^{p+r, q-r+1} \to E_r^{p,q}$ is zero for $r \geq r_0$.
\end{definition}

\begin{proposition}[First quadrant spectral sequences; \ClaimStatusProvedElsewhere]
\label{prop:first-quadrant}
If $E_0^{p,q} = 0$ whenever $p < 0$ or $q < 0$ (first quadrant spectral sequence), 
then the spectral sequence is regular and converges strongly:
\[
E_r^{p,q} \Rightarrow H_{p+q}(C).
\]
\end{proposition}

\begin{proof}
For fixed $(p, q)$ in the first quadrant, the differentials $d_r: E_r^{p,q} \to 
E_r^{p-r, q+r-1}$ must land in the first quadrant. For $r > p$, the target 
$(p-r, q+r-1)$ has $p - r < 0$, so the target is zero. Similarly, 
$d_r: E_r^{p+r, q-r+1} \to E_r^{p,q}$ has source $(p+r, q-r+1)$, which for 
$r > q+1$ has $q - r + 1 < 0$. Thus $d_r = 0$ on and into $E_r^{p,q}$ for 
$r > \max(p, q+1)$.
\end{proof}

\begin{definition}[Strongly convergent]
\label{def:strong-convergence}
A spectral sequence $E_r^{p,q} \Rightarrow H^{p+q}$ is \emph{strongly convergent} 
if:
\begin{enumerate}[label=(\roman*)]
\item The spectral sequence is regular.
\item The filtration on $H^*$ is exhaustive and complete.
\item The isomorphism $E_\infty^{p,q} \cong \gr_p H^{p+q}$ is canonical.
\end{enumerate}
\end{definition}

\begin{theorem}[Zeeman comparison theorem; \ClaimStatusProvedElsewhere]
\label{thm:zeeman}
Let $f: C \to D$ be a filtered chain map between filtered complexes, inducing 
$f_r: E_r(C) \to E_r(D)$ on spectral sequences. If both spectral sequences 
converge strongly and $f_r: E_r(C)^{p,q} \to E_r(D)^{p,q}$ is an isomorphism 
for some $r$ and all $(p, q)$, then:
\[
f_*: H_*(C) \to H_*(D)
\]
is an isomorphism.
\end{theorem}

\begin{proof}
If $f_r$ is an isomorphism, then so is $f_{r+1} = H(f_r)$, and by induction 
$f_\infty: E_\infty(C) \to E_\infty(D)$ is an isomorphism. Strong convergence 
implies $\gr f_*: \gr H_*(C) \to \gr H_*(D)$ is an isomorphism. Completeness 
and exhaustiveness imply $f_*$ itself is an isomorphism.
\end{proof}

\begin{theorem}[Spectral sequence of a filtered dg Lie algebra; \ClaimStatusProvedElsewhere]
\label{thm:spectral-sequence-filtered-dg}
\index{spectral sequence!filtered dg Lie algebra}
Let $(\fg, d, [\cdot,\cdot])$ be a dg Lie algebra with a
complete descending filtration $F^0\fg \supset F^1\fg \supset \cdots$
compatible with the bracket:
$[F^p\fg, F^q\fg] \subset F^{p+q}\fg$.  Then:
\begin{enumerate}[label=\textup{(\roman*)}]
\item There is a spectral sequence $E_r^{p,q} \Rightarrow H^{p+q}(\fg)$
  with $E_0^{p,q} = \gr^p_{p+q}\fg$ and $d_0$ the induced differential on
  the associated graded.
\item The $E_1$ page computes the cohomology of each graded piece:
  $E_1^{p,q} = H^{p+q}(\gr^p\fg)$.
\item The $d_1$ differential is the connecting homomorphism
  in the long exact sequence of the filtration.
\item If the filtration is bounded or the spectral sequence is
  regular, it converges strongly.
\end{enumerate}
\end{theorem}

\begin{proof}
Apply the spectral sequence of a filtered complex
(Theorem~\ref{thm:early-pages} and
Theorem~\ref{thm:classical-convergence}) to the underlying
filtered cochain complex $(C^\bullet(\fg), d)$.
The compatibility $[F^p, F^q] \subset F^{p+q}$ ensures the
filtration is preserved by the Chevalley--Eilenberg
differential.
\end{proof}

\begin{proposition}[Convergence for complete filtrations \cite{Weibel94, Boardman-conditional}; \ClaimStatusProvedElsewhere]
\label{prop:complete-filt-convergence}
Let $(C, F_\bullet)$ be a filtered complex with complete, exhaustive filtration 
bounded below in each degree. The spectral sequence converges to $H_*(C)$ 
(not just to $H_*(\widehat{C})$) if and only if the natural map:
\[
H_*(C) \to \varprojlim_p H_*(C / F_p C)
\]
is an isomorphism (i.e., $\varprojlim^1 H_*(C / F_p C) = 0$).
\end{proposition}


\section{The bar spectral sequence}
\label{sec:bar-spectral-sequence}


\begin{construction}[Bar filtration]
\label{constr:bar-filtration}
For a chiral algebra $\cA$, the geometric bar complex $\Bbar^{\mathrm{geom}}(\cA)$ 
has a filtration by \emph{bar degree}:
\[
F_p \Bbar^{\mathrm{geom}}(\cA) := \bigoplus_{n \leq p} \Bbar^{\mathrm{geom}}_n(\cA)
\]
where $\Bbar^{\mathrm{geom}}_n$ consists of sections over $\FM_n(X)$.

The differential $d = d_{\mathrm{res}} + d_{\mathrm{dR}}$ decomposes as:
\begin{enumerate}[label=(\roman*)]
\item $d_{\mathrm{dR}}: \Bbar^{\mathrm{geom}}_n \to \Bbar^{\mathrm{geom}}_n$ (preserves bar degree).
\item $d_{\mathrm{res}}: \Bbar^{\mathrm{geom}}_n \to \Bbar^{\mathrm{geom}}_{n-1}$ (decreases bar degree by 1).
\end{enumerate}
Thus $d(F_p) \subseteq F_p$, and the filtration is compatible with the differential.
\end{construction}

\begin{theorem}[Bar spectral sequence; \ClaimStatusProvedHere]
\label{thm:bar-ss}
The bar filtration induces a spectral sequence:
\[
E_1^{p,q} = H^{p+q}(\Bbar^{\mathrm{geom}}_p(\cA), d_{\mathrm{dR}}) 
\Longrightarrow H^{p+q}(\Bbar^{\mathrm{geom}}(\cA), d).
\]
The $E_1$-page is the de Rham cohomology of the bar complex at each fixed 
bar degree.
\end{theorem}

\begin{proof}
By Construction~\ref{constr:ass-graded-ss}, the $E_0$-page is:
\[
E_0^{p,q} = \gr_p \Bbar^{\mathrm{geom}}_{p+q} = \Bbar^{\mathrm{geom}}_p(\cA)^{p+q}
\]
(the degree-$(p+q)$ part of the bar-degree-$p$ component). The $d_0$ differential 
is $d_{\mathrm{dR}}$ since this is the component of $d$ preserving filtration 
degree. Thus:
\[
E_1^{p,q} = H^{p+q}(\gr_p \Bbar^{\mathrm{geom}}(\cA), d_{\mathrm{dR}}) 
         = H^{p+q}_{\mathrm{dR}}(\FM_p(X), \cA^{\boxtimes p} \otimes \Omega^\bullet_{\log}).
\]
The $d_1$ differential is induced by $d_{\mathrm{res}}$, the residue part of the 
bar differential.
\end{proof}

\begin{computation}[\texorpdfstring{$E_1$}{E1}-page for Heisenberg]
\label{comp:e1-heisenberg}
For the Heisenberg algebra $\cH$ on $X = \bA^1$.  (We use cohomological notation $H^n_{\mathrm{dR}}$ for de~Rham cohomology entries, matching the main text convention.)

At bar degree $p = 1$:
\[
E_1^{1,q} = H^{1+q}_{\mathrm{dR}}(\bA^1, \cH) = \begin{cases} 
\cH & q = -1 \\ 0 & \text{otherwise} 
\end{cases}
\]
since $H^0_{\mathrm{dR}}(\bA^1) = k$ and higher de Rham cohomology vanishes.

At bar degree $p = 2$:
\[
E_1^{2,q} = H^{2+q}_{\mathrm{dR}}(\FM_2(\bA^1), \cH \boxtimes \cH \otimes \Omega^1_{\log})
\]
The FM compactification $\FM_2(\bA^1) \cong \bA^2$ (the blow-up of $\bA^2$ along the smooth diagonal $\bA^1$ is again isomorphic to~$\bA^2$), with boundary divisor
$D_{12} \cong \bA^1$. The logarithmic de Rham cohomology computes:
\[
E_1^{2,q} = \begin{cases} 
\cH \otimes \cH & q = -1 \\ 0 & \text{otherwise}
\end{cases}
\]
\end{computation}

\begin{proposition}[Degeneration for Koszul algebras; \ClaimStatusProvedHere]
\label{prop:degen-koszul}
If $\cA$ is a Koszul chiral algebra, the bar spectral sequence
degenerates at $E_2$.  The $E_\infty$ page is concentrated on the
Koszul line $q = 0$:
\[
E_\infty^{p,0} \cong (\cA^!)_p, \qquad
E_\infty^{p,q} = 0 \text{ for } q \neq 0,
\]
where $\cA^!$ is the Koszul dual and $(\cA^!)_p$ its weight-$p$ component.
\end{proposition}

\begin{proof}
By Definition~\ref{def:chiral-koszul-pair}, $\cA$ is Koszul if and only if $\mathrm{Ext}^{i,j}_{\cA}(\mathbb{C}, \mathbb{C}) = 0$ for $i \neq j$, where $(i,j)$ is the bigrading by homological and internal degree.

By Theorem~\ref{thm:bar-ss}, the $E_1$ page is $E_1^{p,q} = H^{p+q}_{\mathrm{dR}}(\FM_p(X), \cA^{\boxtimes p} \otimes \Omega^\bullet_{\log})$, with $d_1$ induced by the residue differential. The $E_2$ page is therefore:
\[
E_2^{p,q} = H^q(E_1^{p,*}, d_1) = \mathrm{Ext}^{p,p+q}_{\cA}(\mathbb{C}, \mathbb{C}).
\]
The Koszul condition $\mathrm{Ext}^{i,j} = 0$ for $i \neq j$ translates (setting $i = p$, $j = p + q$) to $E_2^{p,q} = 0$ for $q \neq 0$.

Since $E_2^{p,q} = 0$ for $q \neq 0$, the differentials $d_r: E_r^{p,0} \to E_r^{p-r, r-1}$ for $r \geq 2$ land in bidegree with $q = r - 1 \geq 1 \neq 0$, hence in a zero group.  Similarly, no nonzero differential can arrive at $E_r^{p,0}$.  Thus $d_r = 0$ for all $r \geq 2$, and $E_2 = E_\infty$.

The identification $E_\infty^{p,0} \cong (\cA^!)_p$ follows from $\mathrm{Ext}^{p,p}_{\cA}(\mathbb{C}, \mathbb{C}) = (\cA^!)_p$ by the definition of the Koszul dual.
\end{proof}


\section{The genus spectral sequence}
\label{sec:genus-spectral-sequence}

\medskip\noindent\textbf{Convention change.}  The preceding sections used the homological convention ($d$ lowers total degree, $d_r$ has bidegree $(-r,r{-}1)$).  For the genus spectral sequence we switch to the cohomological convention of the main text: the genus filtration $F^g$ is \emph{decreasing}, and the differential $d_r$ raises the filtration index by~$r$.  The two conventions are related by reindexing $E_r^{p,q} \leftrightarrow E_r^{-p,-q}$ as noted in Construction~\ref{constr:ass-graded-ss}.

\medskip
A filtration by genus organizes quantum corrections.

\begin{construction}[Genus filtration]
\label{constr:genus-filtration}
The following is an algebraic construction; convergence of the genus sum
requires the analytic HS-sewing criterion (Theorem~\ref{thm:general-hs-sewing})
and is not claimed here.  Define the algebraic total bar complex as the direct sum:
\[
\Bbar^{\mathrm{tot}}(\cA) := \bigoplus_{g \geq 0} \Bbar^{(g)}(\cA).
\]
The \emph{genus filtration} is the decreasing filtration by genus:
\[
F^g \Bbar^{\mathrm{tot}}(\cA) := \bigoplus_{h \geq g} \Bbar^{(h)}(\cA).
\]
The total differential $d^{\mathrm{tot}} = d_0 + d_1 + d_2 + \cdots$ decomposes
by genus increase, with $d_k$ raising genus by $k$.
\end{construction}

\begin{theorem}[Genus spectral sequence \cite{BD04}; \ClaimStatusProvedElsewhere]
\label{thm:genus-ss}
The genus filtration induces a spectral sequence:
\[
E_1^{g,n} = H_n(\Bbar^{(g)}(\cA), d_0) \Longrightarrow H_n(\Bbar^{\mathrm{tot}}(\cA))
\]
where:
\begin{enumerate}[label=(\roman*)]
\item $E_1^{0,n}$ is the genus-0 bar homology (classical chiral Hochschild homology).
\item The differential $d_1: E_1^{g,n} \to E_1^{g+1,n}$ encodes one-loop quantum 
corrections.
\item Higher differentials $d_r$ encode $r$-loop corrections.
\end{enumerate}
\end{theorem}

\begin{proposition}[Central charge and \texorpdfstring{$d_1$}{d1}; \ClaimStatusProvedHere]
\label{prop:central-charge-d1}
For a conformal vertex algebra with central charge $c$:
\begin{enumerate}[label=(\roman*)]
\item If $c = 0$, then $d_1 = 0$ and the spectral sequence degenerates at $E_1$.
\item If $c \neq 0$, the $d_1$ differential is nonzero and proportional to $c$.
\item The $E_2$-page encodes the ``one-loop corrected'' chiral homology.
\end{enumerate}
\end{proposition}

\begin{proof}
The $d_1$ differential in the genus spectral sequence (Theorem~\ref{thm:genus-ss}) maps $E_1^{0,n} \to E_1^{1,n}$, encoding the one-loop (genus 1) quantum correction to genus-0 bar homology.

\emph{(i)} If $c = 0$, the genus-1 correction vanishes: the quantum correction to the Arnold relation at genus 1 is proportional to $c \cdot E_2(\tau)$, so $d_1 = 0$ when $c = 0$. With $d_1 = 0$, we have $E_2 = E_1$ and no further differentials can produce nonzero contributions at genus $\geq 2$.  By Theorem~\ref{thm:genus-universality}, all genus-$g$ scalar-lane corrections are proportional to $\kappa(\cA)$; for Virasoro $\kappa(\mathrm{Vir}_c) = c/2$ (AP48: $\kappa = c/2$ is specific to Virasoro; for other families $\kappa$ must be computed from the algebra's own bar complex), so $c = 0$ forces $\kappa = 0$ and all scalar-lane quantum corrections vanish simultaneously.  Thus the spectral sequence degenerates at $E_1$.

\emph{(ii)} If $c \neq 0$, the genus-1 correction $m_0^{(1)} = c \cdot E_2(\tau)$ is nonzero. The $d_1$ differential on the vacuum class is:
\[
d_1([\mathbf{1}]) = c \cdot [\text{genus-1 contribution}] \neq 0,
\]
which is manifestly proportional to $c$.

\emph{(iii)} The $E_2$-page $E_2^{g,n} = H(E_1^{g,n}, d_1)$ incorporates the genus-1 correction into the genus-0 homology, giving the one-loop corrected chiral homology by definition.
\end{proof}

\begin{example}[Heisenberg at higher genus]
\label{ex:heisenberg-genus-ss}
For the Heisenberg algebra $\cH$ with $c = 1$:

At genus 0: $E_1^{0,*} = H_*(\Bbar^{(0)}(\cH)) \cong \cH$ (Koszul).

At genus 1: The differential $d_1: E_1^{0,n} \to E_1^{1,n}$ involves the torus 
partition function. For the vacuum character:
\[
d_1([\mathbf{1}]) \propto \frac{1}{\eta(q)}
\]
where $\eta(q) = q^{1/24} \prod_{n=1}^\infty (1 - q^n)$ is the Dedekind eta function.
\end{example}

\begin{theorem}[Convergence of genus spectral sequence \cite{Weibel94, BD04}; \ClaimStatusProvedElsewhere]
\label{thm:genus-ss-convergence}
\leavevmode\\
The genus spectral sequence converges to $H_*(\Bbar^{\mathrm{tot}}(\cA))$
provided the following conditions hold:
\begin{enumerate}[label=(\roman*)]
\item The genus filtration is complete: $\bigcap_g F^g = 0$.
\item The filtration is bounded below in each homological degree.
\end{enumerate}
For conformal vertex algebras with $c = 0$, the spectral sequence degenerates 
at $E_1$ and:
\[
H_n(\Bbar^{\mathrm{tot}}(\cA)) = H_n(\Bbar^{(0)}(\cA))
\]
(no quantum corrections).
\end{theorem}


\section{The Chevalley--Cousin spectral sequence}
\label{sec:chevalley-cousin}


\begin{construction}[Chevalley--Cousin complex]
\label{constr:chevalley-cousin}
Let $Y$ be a variety with stratification $Y = \bigsqcup_\alpha Y_\alpha$ by 
locally closed subvarieties. For a D-module (or constructible sheaf) $\cM$ on $Y$:

\emph{Step~1.} Order the strata by closure: $Y_\alpha \subseteq \overline{Y_\beta}$ 
implies $\alpha \leq \beta$. Let $Y_{\leq \alpha} = \bigcup_{\beta \leq \alpha} Y_\beta$.

\emph{Step~2.} Define the \emph{Cousin filtration}:
\[
F^p \cM := i_{p,*} i_p^! \cM
\]
where $i_p: Y_{\leq p} \hookrightarrow Y$ is the inclusion of the $p$-th closed 
subset in the stratification ordering.

\emph{Step~3.} The \emph{Chevalley--Cousin complex} is the associated graded:
\[
\gr^p \cM = F^p \cM / F^{p-1} \cM \cong i_{\alpha_p,*} i_{\alpha_p}^! \cM|_{Y_{\alpha_p}}
\]
the $!$-restriction to the $p$-th stratum.
\end{construction}

\begin{theorem}[Chevalley--Cousin spectral sequence \cite{Har77, KS90}; \ClaimStatusProvedElsewhere]
\label{thm:chevalley-cousin-ss}
For a D-module $\cM$ on a stratified variety $Y$:
\[
E_1^{p,q} = H^{p+q}(Y_{\alpha_p}, i_{\alpha_p}^! \cM) \Longrightarrow H^{p+q}(Y, \cM)
\]
where $\alpha_p$ is the stratum of depth $p$.

The $d_1$ differential is the \emph{residue map}:
\[
d_1: H^*(Y_\alpha, i_\alpha^! \cM) \to H^{*+1}(Y_{\beta}, i_\beta^! \cM)
\]
for strata $Y_\alpha$ adjacent to $Y_\beta$ (i.e., $Y_\alpha \subset \partial \overline{Y_\beta}$).
\end{theorem}

\begin{computation}[Chevalley--Cousin for \texorpdfstring{$\FM_3$}{FM3}]
\label{comp:cc-fm3}
For $\FM_3(X)$ with $X = \bA^1$, the stratification is:
\begin{enumerate}[label=(\roman*)]
\item $Y_0 = \Conf_3(\bA^1)$: the open stratum (depth 0).
\item $Y_1 = D_{12} \sqcup D_{13} \sqcup D_{23}$: three boundary divisors (depth 1).
\item $Y_2 = D_{12} \cap D_{23}$: the triple collision (depth 2).
\end{enumerate}

For the trivial D-module $\cO_{\FM_3}$:

$E_1^{0,q} = H^q(\Conf_3(\bA^1), \cO) = \begin{cases} k & q = 0 \\ 0 & q > 0 \end{cases}$

$E_1^{1,q} = H^q(D_{12} \sqcup D_{13} \sqcup D_{23}, i^! \cO) = \begin{cases} k^3 & q = 0 \\ 0 & q > 0 \end{cases}$

$E_1^{2,q} = H^q(D_{123}, i^! \cO) = \begin{cases} k & q = 0 \\ 0 & q > 0 \end{cases}$

The $d_1$ differentials encode how residues along boundary strata relate cohomology 
classes.
\end{computation}

\begin{theorem}[Cousin resolution for holonomic D-modules \cite{KS90}; \ClaimStatusProvedElsewhere]
\label{thm:cousin-resolution}
For a holonomic D-module $\cM$ on a smooth variety $X$, the Chevalley--Cousin 
complex provides a resolution:
\[
0 \to \cM \to i_{0,*} i_0^! \cM \to i_{1,*} i_1^! \cM \to \cdots
\]
where the strata are the connected components of the characteristic variety 
$\mathrm{Ch}(\cM) \subset T^*X$.

For D-modules with regular singularities, this spectral sequence degenerates 
at $E_2$ and:
\[
H^*(X, \cM) \cong E_2^{*,0}.
\]
\end{theorem}


\section{Multiplicative spectral sequences}
\label{sec:mult-ss}


\begin{definition}[Multiplicative spectral sequence]
\label{def:mult-ss}
A spectral sequence $(E_r^{*,*}, d_r)$ is \emph{multiplicative} if:
\begin{enumerate}[label=(\roman*)]
\item Each $E_r^{*,*}$ is a bigraded algebra.
\item The differential $d_r$ is a derivation: $d_r(xy) = (d_r x) y + (-1)^{|x|} x (d_r y)$.
\item The product on $E_{r+1} = H(E_r, d_r)$ is induced from $E_r$.
\end{enumerate}
\end{definition}

\begin{proposition}[Bar spectral sequence is multiplicative \cite{LV12}; \ClaimStatusProvedElsewhere]
\label{prop:bar-ss-mult}
The bar spectral sequence for an algebra $A$ is multiplicative with product:
\[
[a_1|\cdots|a_m] \cdot [b_1|\cdots|b_n] = \sum_{\sigma \in \mathrm{Sh}(m,n)} 
\epsilon(\sigma) [c_{\sigma(1)}|\cdots|c_{\sigma(m+n)}]
\]
where $\{c_1, \ldots, c_{m+n}\} = \{a_1, \ldots, a_m, b_1, \ldots, b_n\}$ and 
$\mathrm{Sh}(m,n)$ is the set of $(m,n)$-shuffles.

The shuffle product gives $\Bbar(A)$ an algebra structure; together with the deconcatenation comultiplication, $\Bbar(A)$ becomes a differential graded bialgebra.
\end{proposition}

\begin{theorem}[Convergence of multiplicative spectral sequences \cite{Weibel94}; \ClaimStatusProvedElsewhere]
\label{thm:mult-ss-conv}
If a multiplicative spectral sequence $E_r \Rightarrow H$ converges strongly, 
then:
\begin{enumerate}[label=(\roman*)]
\item $H$ inherits an algebra structure.
\item The filtration on $H$ is compatible with the product.
\item $\gr H \cong E_\infty$ as algebras.
\end{enumerate}
\end{theorem}

\begin{example}[Adams spectral sequence]
\label{ex:adams-ss}
The Adams spectral sequence in stable homotopy theory is multiplicative:
\[
E_2^{s,t} = \Ext_{\cA}^{s,t}(\bF_p, \bF_p) \Longrightarrow \pi_{t-s}(S^0)_{(p)}
\]
where $\cA$ is the Steenrod algebra. The product on $E_2$ is the Yoneda product 
in $\Ext$, and the product on $E_\infty$ is induced from the composition product 
in stable homotopy groups.

The connection to chiral algebras runs through factorization
homology and stable homotopy theory.
\end{example}
